\section{Two distinct stability mechanisms}\label{app:separation}
The curvature and same-report alternatives in the informative Bayesian theorem~\cite{GTF32} are genuinely different. The examples below address that comparison directly. They do not assert that a stability certificate alone proves the additional geometric hypotheses.

For a homogeneous affine Bayesian experiment on a fixed compact convex domain $K$, write
\[
 \Phi(u)=\int\norm{F(y,u)}^2\ell(y,u)\,d\lambda(y).
\]
The curvature route asks for $\frac12D^2\Phi(u)[v,v]\le b_t\norm v^2$ with products of the $b_t$ decaying uniformly. The same-report route asks for uniformly bounded metrics dominating the task norm and deterministic bounds on $\operatorname{dist}_t(F(y,u),F(y,v))$ whose products decay uniformly for all common nonexceptional reports. The latter is not merely an average contraction condition.

\begin{proposition}[Neither stability route implies the other]\label{prop:separation}
There are explicit positive prepared finite-state experiments with the following properties.

\textup{(a)} A two-state Gaussian experiment satisfies the same-report route on its full invariant prediction domain, but every valid uniform Euclidean curvature factor is greater than $48020000/46268279>1$.

\textup{(b)} A two-state hold--reset experiment has curvature factor $a$ for any chosen $0<a<1$, but cannot satisfy the same-report route with uniformly bounded, norm-dominating metrics. This second example separates stability certificates; its atomic row laws do not satisfy the positive-dimensional small-ball premise of the older regular-geometry theorem.
\end{proposition}
\begin{proof}
For (a), let $q=1/100$, $\lambda=1-2q=49/50$, and
\[
 P=\begin{pmatrix}1-q&q\\q&1-q\end{pmatrix},\quad
 f_\pm(y)=\varphi(y\mp1/2),\quad L_y=P\operatorname{diag}(f_+(y),f_-(y)).
\]
This example uses \emph{report first, transition second}, exactly as in the older informative filtering certificate; it is separate from Section~\ref{sec:gaussian}. The invariant domain is $K=P\Delta_1$, represented by $p\in[q,1-q]$. Its categorical forecast is $(p,1-p)$. Gaussian densities are positive and every interior point of $K$ is in the support of every acquired posterior row.

In the log-odds metric $d(p,p')=|\log(p/(1-p))-\log(p'/(1-p'))|$, a Bayesian report adds the same log likelihood ratio and is an isometry. The transition $p\mapsto q+\lambda p$ has derivative in log-odds coordinates
\[
 \frac{\lambda p(1-p)}{(q+\lambda p)(1-q-\lambda p)}\le\lambda<1,
\]
since the denominator minus $p(1-p)$ equals $q(1-q)(2p-1)^2\ge0$. Thus the full normalized update contracts by $\lambda$. On $K$ this metric has finite diameter and is uniformly equivalent to the task norm. It supplies the same-report route.

For the curvature route, direct differentiation of the raw perspective gives, relative to the categorical tangent norm $\norm{(v,-v)}^2=2v^2$, the exact factor
\[
 \beta(p)=\lambda^2\int
       \frac{f_+(y)^2f_-(y)^2}
            {\{p f_+(y)+(1-p)f_-(y)\}^3}\,dy.
\]
On a compact interior neighborhood of the point used next the integrands and derivatives are bounded by a constant times $f_++f_-$, justifying differentiation under the integral. Since $f_+/f_-=e^y$ and $f_+^2/f_-=e\,\varphi(y-3/2)$, restriction to $y\le2$ gives
\[
 \beta(p)\ge
 \frac{\lambda^2e\,\Phi(1/2)}{\{1+p(e^2-1)\}^3}
 >\frac{\lambda^2(27/10)(1/2)}{(1+7p)^3}.
\]
Here $e>27/10$, $e^2<8$ follow directly from the exponential series, and $\Phi(1/2)>1/2$ by symmetry. The point $p=11/1000$ is strictly inside $K$. Substitution gives $48020000/46268279>1$. Therefore every homogeneous-domain curvature factor must exceed this number at every phase, and its products cannot decay. This is a failed sufficient curvature certificate, not a failed filtering theorem.

For (b), take hidden categorical states with prediction $u\in\Delta_1$, prepare the known law $u_0=(3/4,1/4)^T$, and use two report symbols $H,R$ with raw instruments
\[
 L_H=aI,\qquad L_R=(1-a)\pi\one^T,
 \qquad \pi=\tfrac12\binom11.
\]
Report $H$ keeps the hidden state, while report $R$ resets it to $\pi$. Both probabilities are positive and independent of the input. The normalized updates are $F_H(u)=u$ and $F_R(u)=\pi$, so
\[
 \Phi(u)=a\norm u^2+(1-a)\norm\pi^2,
 \qquad\tfrac12D^2\Phi(u)[v,v]=a\norm v^2.
\]
Thus the curvature route holds with geometric factor $a<1$. At every positive time the actually acquired prediction law assigns positive mass to both $u_0$ (no reset yet) and $\pi$ (at least one reset). The counterexample therefore persists on their convex hull, and does not rely on states that the preparation never acquires.

If the same-report route held, compose its inequalities along the report word $H^n$, which has positive probability for every $n$. For any two distinct $u,v$ in the declared domain this would give
\[
 \operatorname{dist}_n(u,v)\le C_0\rho^n\operatorname{dist}_0(u,v).
\]
Uniform domination of $\norm{u-v}$ by the left side and bounded initial diameter contradict this as $n\to\infty$. A zero-probability exception cannot remove the positive-mass symbol $H$. This proves the asserted separation.
\end{proof}
